\section{An improved exponent without an epsilon loss}\label{ep:section}

This section is new in version 2. The original power bound and the
exceptional-family proof above are retained with their original numbering.
Write $s=\log_3 2$, $\varphi=(1+\sqrt5)/2$, and $g=\log_3\varphi$.

\begin{theorem}\label{ep:main}
There are absolute constants $c_*>0$ and $X_0$ such that
\[
 \MC(X)\ge c_*X^{\beta_*}\quad(X\ge X_0),\qquad
 \beta_*={ (1+s)^2\over7-5s}=0.691726633148796\ldots.
\]
The same lower bound holds for $\MH(X)$. No subtraction of an arbitrary
positive epsilon from $\beta_*$ is required.
\end{theorem}

\subsection{A boundary estimate for the mean Fourier discrepancy}
Use $F_l,f_l,\D_l$ from Section 2. Their second and fourth moments are
$2^{-l}$ and $(3/8)^l$. Let $T_D$ be all nonzero reduced fractions on
$\T$ with denominator at most $D$ coprime to 3, and put
$S_D(L)=\sum_{z\in T_D}f_L(z)$ and $h=\lceil2\log_3D\rceil$.
Multiplication by 3 permutes $T_D$, whose separation is at least $D^{-2}$.
The sampling lemma gives, uniformly for $0\le m\le h$,
\[
 \sum_{T_D}f_m^2\ll D^2 2^{-m},\qquad
 \sum_{T_D}f_m^4\ll D^2(3/8)^m.
\]
Indeed $3^h<3D^2$, so the degrees in both applications are $O(D^2)$.
Set
\[
 a_*={3s-1\over4},\quad A_*={2(3-s)\over3s-1},\quad
 \eta_0=3-4s,\quad r_*={3-s\over2}=a_*A_*.
\]
Here $4<A_*<6$, $0<\eta_0<s$, and $r_*>1$.
For $2h\le L<3h$ write $l=L-2h$ and apply Holder's inequality
with exponents $2,4,4$ to blocks of lengths $h,h,l$. This yields
\begin{equation}\label{ep:farey}
 S_D(L)\ll D^2 2^{-h/2}(3/8)^{(h+l)/4}
 \le C D^{(3+s)/2}3^{-a_*L}.
\end{equation}
The last inequality uses $s-1<0$ and $h\ge2\log_3D$.
For $L\ge3h$, dropping extra factors instead gives
\begin{equation}\label{ep:long}
 S_D(L)\le S_D(3h)\ll D^{\eta_0}.
\end{equation}
For each fixed $D$, also $S_D(L)\to0$: every nonzero fraction has a
finite orbit under multiplication by 3, with all factors
$|(1+\exp(2\pi iz))/2|<1$.

Let $S_Y=\{N\le Y:N\in\D,\ N\equiv1\pmod3\}$ and
$R_N(L)=\sum_{u=1}^{N-1}f_L(u/N)$.
\begin{lemma}\label{ep:mean}
There is an absolute $C_*$ such that, for every fixed $H\ge0$ and
every integer choice $L=L(Y)\ge A_*\log_3Y+H$,
\[
 \limsup_{Y\to\infty}Y^{-s}\sum_{N\in S_Y}R_N(L)
 \le C_*3^{-a_*H}.
\]
\end{lemma}
\begin{proof}
Fixing the high ternary digits shows that
$T_Y(d)=\#\{N\in S_Y:d\mid N\}\le4(Y/d)^s$ for $1\le d\le Y$:
choose a low block with $3^j\le d<3^{j+1}$; its values lie in an
interval of length less than $d$, hence contribute at most one multiple
for each high block. Reducing $u/N$ gives the exact identity
\[
 \sum_{N\in S_Y}R_N(L)=
 \sum_{\substack{2\le d\le Y\\(d,3)=1}}T_Y(d)
       \sum_{\substack{1\le u<d\\(u,d)=1}}f_L(u/d).
\]
In a fixed dyadic range $D/2<d\le D$, its contribution divided by
$Y^s$ is $O(D^{-s}S_D(L))$. One may take $D$ to be powers of 2,
with the final partial range bounded by $D<2Y$.
For ranges with $3h\le L$, equation~\eqref{ep:long} supplies the
summable majorant $CD^{\eta_0-s}$. Each fixed range tends to zero,
so dominated convergence makes their total $o(1)$.
For all remaining ranges, $2h\le L$ once $Y$ is sufficiently large,
because $A_*>4$. Equation~\eqref{ep:farey} gives
\[
 D^{-s}S_D(L)\ll D^{r_*}3^{-a_*L}
 \le 3^{-a_*H}(D/Y)^{r_*}.
\]
The sum of $(D/Y)^{r_*}$ over these dyadic ranges is bounded absolutely.
This proves the stated limsup. In particular the estimate asserts an
arbitrarily small fixed proportion, not necessarily $o(Y^s)$ at fixed $H$.
\end{proof}

\subsection{Affine tails and a fixed residue class}
\begin{lemma}\label{ep:affine}
If $r\in1+3\Z_3$, $r\ne1$, $e=v_3(r-1)$, and $\xi\in\Z_3$, then
\[
 \#\{z\in\D_m:[rz+\xi]_m\in\D_m\}
 \le H_m(e)\le\varphi^m(2/\varphi)^e,
\]
where $[x]_m$ is the residue in $[0,3^m)$ and $H_m(e)$ counts the
independent sets of the graph on $0,\ldots,m-1$ with edges of length $e$.
\end{lemma}
\begin{proof}
Write $r=1+3^eu$, with $u$ a unit. In the addition
$z+3^euz+\xi$, let $B_j$ be the sum of the incoming carry,
the digit of $\xi$, and the shifted digit of $uz$.
On an admissible path the outgoing carry is
$\lfloor(B_j+1)/3\rfloor$, independent of the admissible choice of
$z_j$. Residues $B_j=0,1,2\pmod3$ allow, respectively, two choices,
only $z_j=0$, or only $z_j=1$.
At position $i+e$, after fixing $z_{<i}$, the output digit is
$z_{i+e}+u_0z_i+Q_i(z_{<i})\pmod3$, with $u_0\in\{1,2\}$.
Thus one of the two choices at $i$ forces the digit at $i+e$,
independently of intermediate admissible digits. Backward induction
on the remaining positions bounds the tree by the independent-set
recursion: at a free vertex one either removes that vertex, or removes
it together with its future neighbor; an already forced vertex has
only one branch. Discarding further restrictions only enlarges the
bound. The graph splits into $e$ chains, including empty chains.
A chain of length $l$ has $F_{l+2}\le2\varphi^{l-1}$ independent sets,
also for $l=0$. Multiplying proves the second inequality.
\end{proof}

Fix $c=112=(11011)_3$ and $d=128$ throughout the rest of the proof.
Define
\[
 C_d(c)=\#\{a,b,a',b'\in\D_d:
              c(a-a')+(b-b')\equiv0\pmod{3^d}\}.
\]
For every right-hand side the corresponding count is at most $C_d(c)$,
by Cauchy--Schwarz applied to the autocorrelation of
$h(z)=\#\{(a,b)\in\D_d^2:ca+b\equiv z\pmod{3^d}\}$.
Also $C_d(c)\ge(16/3)^d$.

The following finite certificate specifies the constant used here
without a numerical spectral-radius approximation. On states
$-56,\ldots,56$ put $T_{u,v}=w(3v-u)$, where
\[
 w(0)=4,\quad w(\pm1)=2,\quad w(\pm111)=1,\quad
 w(\pm112)=2,\quad w(\pm113)=1
\]
and all other weights vanish. These are precisely the digit-difference
multiplicities. The state interval is closed under every integral
transition. Starting with $v_0(u)=\mathbf1_{u=0}$, iterate
$v_{j+1}(v)=\sum_u v_j(u)T_{u,v}$. Then
$C=C_d(c)=\sum_u v_{128}(u)$, with unrestricted final carry, equals
\[
\begin{split}
 C={}&2848039325767103655925435712986932980549021547392299357785\\[-2pt]
     &\hspace{8mm}070566911039704321748559183306903592
\end{split}
\]
where the two lines are concatenated decimal digits.
The exact integer checks supplied with the article give
\begin{equation}\label{ep:cert}
 C^{1161}3^{148608}\varphi^{491392}<2^{937216},\qquad
 2^{20805}<3^{13127}.
\end{equation}
These comparisons require integers only: for $m=491392$, write
$\varphi^m=F_m\varphi+F_{m-1}$, multiply by
$C^{1161}3^{148608}$ to obtain $A\varphi+B$, and check
$D=2^{937217}-A-2B>0$ and $D^2>5A^2$.
The executable certificate recomputes the matrix both forwards and
backwards and performs these comparisons exactly.
Set $\lambda=C^{1/d}$ and $\tau=\log_3\lambda$.

\begin{lemma}\label{ep:pairs}
Let $Y=3^k$, $k>d$, and $G\subset\{N\in\D_k:N\equiv1\pmod{3^d}\}$.
For $Y<X$ put
$A_N(X)=\{1\le M\le X:M\equiv c\pmod{3^d},\ NM\in\D\}$.
Then
\[
 \sum_{N<N'\in G}|A_N(X)\cap A_{N'}(X)|\ll_d X^gY^{\tau-g}.
\]
\end{lemma}
\begin{proof}
For four prefixes $N',x,N,y$ of length $m\ge d$ with low blocks
$1,c,1,c$ and $N'x\equiv Ny\pmod{3^m}$, adding four blocks of
length $d$ reduces the next congruence to
\[
 R+c(a-a')+(b-b')\equiv0\pmod{3^d},\qquad R=(N'x-Ny)/3^m.
\]
The quadratic terms vanish since $m\ge d$. There are at most $C$
extensions. If the first distinct denominator digit occurs at $e$,
there are at most $4^{e-d}$ common prefixes and two choices at $e$.
Applying full blocks and bounding a remaining partial block trivially
by $16^d$ proves that the number of length-$k$ quadruples is
$O_d(4^e\lambda^{k-e})$.

Choose $n\ge k$ with $X<3^n\le3X$. For each pair, the map
$M\mapsto[NM]_n$ is injective, since $N$ is a unit. Conditional on a
prefix $x_0$ of length $k$, its remaining digits satisfy an affine
condition with ratio $r=N'/N$ and shift
$\xi=(rx_0-[rx_0]_k)/3^k$. Lemma~\ref{ep:affine} bounds these tails.
Summing over $e$ therefore gives
\[
 O_d\left(\varphi^{n-k}\lambda^k
       \sum_{e\ge d}\left({8\over\varphi\lambda}\right)^e\right).
\]
The ratio is at most $3/(2\varphi)<1$, since $\lambda\ge16/3$.
This proves the assertion, including arbitrarily close distinct pairs.
\end{proof}

\subsection{The first moment and completion of the proof}
Put
\[
 \alpha_*={1\over A_*-1}={3s-1\over7-5s},\qquad
 Y=3^k=3^{\lfloor\log_3(c_0X^{\alpha_*})\rfloor},
\]
where $c_0>0$ will be fixed before $X\to\infty$.
The shell
$\mathcal S_{k,d}=\{N\in\D_k:N\equiv1\pmod{3^d},\ N\ge3^{k-1}\}$
has exactly $2^{-d-1}Y^s$ members, all in $[Y/3,Y/2)$.
Choose a fixed $H$ with $C_*3^{-a_*H}<2^{-d-5}$, and then
$c_0\le3^{-(H+d+4)/(A_*-1)}$. For
$L_* =\lfloor\log_3(XY/3)\rfloor-d$ one has
$L_*-A_*\log_3Y\ge H+2$.
Lemma~\ref{ep:mean} and Markov's inequality show that at least half
the shell, denoted $G$, has $R_N(L_*)\le1/2$ for all sufficiently large $X$.

For $L_N=\lfloor\log_3(NX)\rfloor$, count multiples of $N$ of the form
$z=c+3^dw$, $w\in\D_{L_N-d}$. Orthogonality, permutation of residues
by $3^d$, and monotonicity of $f_l$ give at least
\[
 {2^{L_N-d}\over N}(1-R_N(L_N-d))
 \ge2^{-d-2}X^sN^{s-1}
\]
such multiples. They are positive and less than $NX$, and $z/N$ lies
in the required residue class. Consequently
\begin{equation}\label{ep:first}
 \sum_{N\in G}|A_N(X)|\gg_d X^sY^{2s-1}.
\end{equation}
The exact inequalities~\eqref{ep:cert} imply
$\alpha_*<\alpha_0=1161/5000$ and
$g-s+\alpha_0(1+\tau-g-2s)<0$.
The coefficient $1+\tau-g-2s$ is positive, since
$\tau\ge4s-1$ and $2s>g$. Hence
\[
 g-s+\alpha_*(1+\tau-g-2s)<0.
\]
Lemma~\ref{ep:pairs} is thus $o(1)$ times~\eqref{ep:first}.
Bonferroni's inequality gives
\[
 \MC(X)\ge\left|\bigcup_{N\in G}A_N(X)\right|
 \gg X^sY^{2s-1}\gg X^{s+(2s-1)\alpha_*}=X^{\beta_*}.
\]
The positive-witness inclusion proved above transfers the bound to
$\MH$. This completes the proof of Theorem~\ref{ep:main}.
The constants can be very small and the threshold very large; no
practical search threshold or sharp asymptotic is asserted.
